\chapter{Introducción}

\begin{itemize}
    \item Contextualización del problema
    \item Relevancia del estudio 
    \item Propósito del trabajo 
    \item Estructura del documento 
\end{itemize}

\section{Planteamiento del problema}
La introducción presenta el problema que aborda la investigación en su contexto de aplicación, y la solución propuesta en el trabajo. Incluye referencias de los trabajos previos realizados en la misma temática (estado del arte), y debe resaltar la diferencia entre esos trabajos y la metodología propuesta. Es importante que queden claros los aspectos innovadores, aportes y mejoras logradas con la metodología propuesta. Para mayor información sobre cómo redactar el documento de tesis, consultar Katz \cite{abramson2015writing}.


\section{Justificación}
Aquí además de una breve justificación acerca de por qué el trabajo es importante a nivel académico, social, etc, también describe los objetivos específicos y el objetivo general de manera resumida y en párrafo, procura no usar viñetas ni nada por el estilo, sólo redacción.

\section{Estado del arte, antecedentes y conceptos claves}
Revisa los aspectos conceptuales básicos necesarios para comprender el problema y la solución propuesta.




